%%%%%%%%%%%%%%%%%%%%%%%%%%%%%%%%%%%%%%%%%
% University Assignment Title Page 
% LaTeX Template
% Version 1.0 (27/12/12)
%
% This template has been downloaded from:
% http://www.LaTeXTemplates.com
%
% Original author:
% WikiBooks (http://en.wikibooks.org/wiki/LaTeX/Title_Creation)
%
% License:
% CC BY-NC-SA 3.0 (http://creativecommons.org/licenses/by-nc-sa/3.0/)
% 
% Instructions for using this template:
% This title page is capable of being compiled as is. This is not useful for 
% including it in another document. To do this, you have two options: 
%
% 1) Copy/paste everything between \begin{document} and \end{document} 
% starting at \begin{titlepage} and paste this into another LaTeX file where you 
% want your title page.
% OR
% 2) Remove everything outside the \begin{titlepage} and \end{titlepage} and 
% move this file to the same directory as the LaTeX file you wish to add it to. 
% Then add \input{./title_page_1.tex} to your LaTeX file where you want your
% title page.
%
%%%%%%%%%%%%%%%%%%%%%%%%%%%%%%%%%%%%%%%%%
%\title{Title page with logo}
%----------------------------------------------------------------------------------------
%	PACKAGES AND OTHER DOCUMENT CONFIGURATIONS
%----------------------------------------------------------------------------------------

\documentclass[12pt]{article}
\usepackage[italian]{babel}
\usepackage[utf8x]{inputenc}
%\usepackage{amsmath}
\usepackage{graphicx}
%\usepackage[colorinlistoftodos]{todonotes}
\usepackage{float}

\usepackage[hidelinks]{hyperref}
\usepackage[all]{hypcap}

\usepackage[table]{xcolor}

\begin{document}
\pagenumbering{roman}

\begin{titlepage}

\newcommand{\HRule}{\rule{\linewidth}{0.5mm}} % Defines a new command for the horizontal lines, change thickness here

\center % Center everything on the page
 
%----------------------------------------------------------------------------------------
%	HEADING SECTIONS
%----------------------------------------------------------------------------------------

\textsc{\LARGE Università degli studi di Milano-Bicocca}\\[1cm] % Name of your university/college
\textsc{\Large Advanced Machine Learning }\\[0.3cm] % Major heading such as course name
\textsc{\large Final Project}\\[0.1cm] % Minor heading such as course title

%----------------------------------------------------------------------------------------
%	TITLE SECTION
%----------------------------------------------------------------------------------------

\HRule \\[0.4cm]
{ \huge \bfseries Mercari Price Suggestion Challenge}\\[0.4cm] % Title of your document
\HRule \\[1.5cm]
 
%----------------------------------------------------------------------------------------
%	AUTHOR SECTION
%----------------------------------------------------------------------------------------

\large
\emph{Authors:}\\
Gabriele Ferrario -  817518 - g.ferrario@campus.unimib.it \\   % Your name
Riccardo Pozzi - 807857 - r.pozzi@campus.unimib.it   \\[1cm] % Your name

% If you don't want a supervisor, uncomment the two lines below and remove the section above
%\Large \emph{Author:}\\
%John \textsc{Smith}\\[3cm] % Your name

%----------------------------------------------------------------------------------------
%	DATE SECTION
%----------------------------------------------------------------------------------------

{\large \today}\\[2cm] % Date, change the \today to a set date if you want to be precise

%----------------------------------------------------------------------------------------
%	LOGO SECTION
%----------------------------------------------------------------------------------------

\includegraphics{logo.png}\\[1cm] % Include a department/university logo - this will require the graphicx package
 
%----------------------------------------------------------------------------------------

\vfill % Fill the rest of the page with whitespace

\end{titlepage}

\pagenumbering{arabic}
\begin{abstract}
    % The ABSTRACT is not a part of the body of the report itself. Rather, the
% abstract is a brief summary of the report contents that is often separately
% circulated so potential readers can decide whether to read the report. The
% abstract should very concisely summarize the whole report: why it was written,
% what was discovered or developed, and what is claimed to be the significance of
% the effort. The abstract does not include figures or tables, and only the most
% significant numerical values or results should be given.

Il progetto si incentra su di un problema di regressione derivato dalla Kaggle
Challenge Mercari Price Suggestion che consiste nella predizione del prezzo di
vendita dei prodotti a partire da feature testuali e non. È stata effettuata una
valutazione sia in termini di performance di regressione che di costi
computazionali di diversi approcci al problema, soprattutto relativi alla
rappresentazione del testo, tra cui Bag of Words e Word Embedding. È stato
valutato il contributo del pre-processing del testo e l'efficacia di diversi
tipi di layer neurali nella gestione del testo. In conclusione sono state
ottenute ottime performance utilizzando un Word Embedding allenabile tramite il
layer Embedding di Keras in combinazione a layer LSTM bidirezionali. Bag of
Words è risultata una buona alternativa pur essendo più semplice e
computazionalmente più economico. Per quanto riguarda il pre-processing del
testo, l'approccio condotto non ha portato miglioramenti rilevanti.
\end{abstract}

\section{Introduction}
% The introduction should provide a clear statement of the problem posed by the
% project, and why the problem is of interest. It should reflect the scenario, if
% available. If needed, the introduction also needs to present background
% information so that the reader can understand the significance of the problem. A
% brief summary of the hypotheses and the approach your group used to solve the
% problem should be given, possibly also including a concise introduction to
% theory or concepts used later to analyze and to discuss the results.


% Descrizione della challenge originale: mercari price. Descrizione di cosa la
% challenge implica:
% - regressione a partire da features categoriche e testuali Dire che le
%   testuali sono importanti e più complicate (anche computazionalmente) e dire
%   che procediamo con un confronto dei diversi approcci esistenti per trattare
%   testi valutandoli in base alle performance ottenute e alla risorse
%   computazionali richieste però dei vari dati se ne parla nella sezione
%   dataset quindi forse è troppo presto
%   per parlare di testo


Il progetto trae origine dalla \textit{Kaggle challenge} \textbf{Mercari Price
Suggestion Challenge} \cite{mercari-price-suggestion-challenge} aperta a fine
Novembre 2017 che come viene reso chiaro dal sottotitolo \textit{"Can you
automatically suggest product prices to online sellers?"} pone l'obiettivo di
stimare il prezzo dei prodotti a partire da alcune loro caratteristiche. \\
Alla base di ciò vi è l'esigenza dell'e-commerce \textit{Mercari} \cite{mercari}
di offrire ai propri venditori un suggerimento sul prezzo di vendita dei
prodotti inseriti.

Si tratta quindi di un problema di \textit{regressione} che a partire dalle
varie caratteristiche dei prodotti, testuali e non, vuole calcolare il prezzo da
suggerire.
\\
Durante lo svolgimento del progetto si valuteranno vari approcci al problema,
sopratutto per quanto riguarda i dati di tipo testuale, sia in termini di errore
rispetto ai dati di \textit{train} che di costi computazionali.

In particolare saranno valutati i seguenti modelli di rappresentazione del
testo: Bag of Words, Tf-Idf, Word Embedding, Word Embedding pre-allenato (GloVe),
Feature extraction con Transfomer pre-allenato (DistilBert); verrà inoltre
valutata l'efficacia del pre-processing del testo.



\section{Datasets}
In this section the available data sets must be presented. The term dataset refers to any type of information source, for example web services for geolocation fall into this category. 
In addition, all necessary data manipulation processes, such as cleaning and enrichment with external sources, must be presented and discussed.



Il Dataset consiste in un elenco di 1391082 prodotti descritti tramite le seguenti caratteristiche:
\subsubsection{Price}
textbf{Price} rappresenta il prezzo per il quale l'articolo è stato venduto (variabile target).
Il prezzo medio è di circa \$26, con un valore minimo pari a \$0 e un valore massimo pari a \$2009; inoltre, presenta una deviazione standard di circa \$38.
Analizzando i percentili ci si accorge che i prezzi sono relativamente bassi in quanto il 75\% dei prodotti hanno un prezzo al di sotto di \$29.
\subsubsection{Shipping}
\textbf{Shipping} è caratterizzato dal valore 1 se la tassa di spedizione è a carico del venditore, altrimenti 0 se è a carico dell'acquirente.
Questo attributo è decentemente ripartito tra i venditori e gli acquirenti, in quanto il 55\% dei prodotti prevede un valore di 0.
Analizzando i prezzo degli articoli ci si aspetta che per quelli che la tassa di spedizione è a carico del venditore avranno un prezzo più alto. Tuttavia, ci sono una serie di fattori contrastanti. Questo può essere vero all'interno di specifiche categorie di prodotti e condizioni degli articoli, ma non quando si confrontano gli articoli sul totale. 
Infatti, il prezzo medio pagato dagli utenti che devono pagare le spese di spedizione (circa \$30) è superiore a quelli che non richiedono costi di spedizione aggiuntivi (circa \$22).
\subsubsection{Category Name}
\textbf{Category\_name} rappresenta la categoria di prodotto a cui appartiene l'articolo.
Nel dataset sono presenti 1287 categorie univoche e tra ognuna di esse si vede una categoria principale/generale, seguita da due o più sottocategorie più specifiche (ad esempio: Women/Tops \& Blouses/T-Shirts). Inoltre, ci sono 6327 articoli che non hanno una categoria assegnata.
Infine, analizzando le dieci categorie più popolari, si nota che l'abbigliamento femminile è molto popolare su Mercari. Infatti, di queste prime dieci categorie 5 sono di abbigliamento femminile; Anche il trucco e l'elettronica sono categorie molto quotate.
\subsubsection{Brand Name}
\textbf{Brand\_name} rappresenta il marchio dell'articolo; nel dataset sono presenti 4809 valori differenti e 632682 valori mancanti.
\subsubsection{Item Description}
textbf{Item\_description} rappresenta la descrizione del prodotto sotto forma di dato non strutturato. Nel dataset sono presenti 4 istanze senza descrizione e 82494 descrizioni con la stringa "no description yet".
Inoltre, non esiste una correlazione tra la lunghezza delle descrizioni e il prezzo, in quanto c'è una correlazione di 0.048.
\textbf{AGGIUNGERE WORDCLOUD}

\section{The Methodological Approach}
% This is the central and most important section of the report. Its objective must
% be to show, with linearity and clarity, the steps that have led to the
% definition of a decision model. The description of the working hypotheses,
% confirmed or denied, can be found in this section together with the description
% of the subsequent refining processes of the models. Comparisons between
% different models (e.g. heuristics vs. optimal models) in terms of quality of
% solutions, their explainability and execution times are welcome.

% Do not attempt to describe all the code in the system, and do not include large
% pieces of code in this section, use pseudo-code where necessary. Complete source
% code should be provided separately (in Appendixes, as separated material or as a
% link to an on-line repo). Instead pick out and describe just the pieces of code
% which, for example:

% \begin{itemize}
%     \item are especially critical to the operation of the system;
%     \item you feel might be of particular interest to the reader for some reason;
%     \item  illustrate a non-standard or innovative way of implementing an algorithm, data
%           structure, etc..
% \end{itemize}

% You should also mention any unforeseen problems you encountered when implementing the
% system and how and to what extent you overcame them. Common problems are:
% difficulties involving existing software.

% todo meglio metterla qui l'analisi del dataset? media, std,... usate dalla normale
\subsubsection{Training e valutazione}

\subsubsection{Dataset split}

Al fine di valutare i vari modelli più correttamente possibile il dataset è
stato suddiviso in \textit{training-set}, \textit{validation-set} e
\textit{test-set}. Quindi il training è stato usato per l'allenamento, il
validation per accertarsi che il modello non sia propenso all'overfitting ed il
test per la valutazione finale.

\subsubsection{Valutazione dell'errore}
% loss e metriche

Per quanto riguarda le performance in termini di errore la funzione di
\textit{loss} utilizzata in fase di training è il \textit{Root Mean Squared
Logaritmic Error (RMSLE)}, in quanto è la misura scelta dalla Kaggle challenge
per confrontare le performance dei vari partecipanti. Inoltre essa risulta
adeguata al problema considerando il vasto intervallo dei valori dei prezzi.

Di seguito la definizione e alcune osservazioni su di RMSLE insieme alle
ulteriori metriche utilizzate; ad esempio il \textit{Mean Absolute Error (MAE)},
che fornisce una più immediata comprensione rispetto al RMSLE. \\
Mean Absolute Error (MAE):
\begin{equation}
    \frac{1}{n} \sum_{i=1}^{n} | y_i - \hat{y_i} |
\end{equation}
\\
Mean Squared Error (MSE):
\begin{equation}
    \frac{1}{n} \sum_{i=1}^{n} ( y_i - \hat{y_i} )^2
\end{equation}
Il quadrato fornisce un peso maggiore agli errori dei valori più elevati\cite{rmse-or-mae}.
\\
Root Mean Squared Error (RMSE):
\begin{equation}
    \sqrt{ \frac{1}{n} \sum_{i=1}^{n} ( y_i - \hat{y_i} )^2}
\end{equation}
La radice dell'MSE, nella stessa scala dei valori.
\\
Mean Squared Logarithmic Error (MSLE):
\begin{equation}
    \frac{1}{n}
        \sum_{i=1}^{n}
            ( \log(y_i+1) - \log(\hat{y_i}+1) )^2
    =
    \frac{1}{n}
        \sum_{i=1}^{n}
            \log^2(\frac{y_i+1}{\hat{y_i}+1})
\end{equation}
Essendo calcolato a partire da un rapporto riflette l'errore relativo e di
conseguenza risulta efficace laddove i valore assoluti presentano variazioni considerevoli.
\\
Root Mean Squared Logarithmic Error (RMSLE):
\begin{equation}
    \sqrt{ 
        \frac{1}{n}
            \sum_{i=1}^{n}
                ( \log(y_i+1) - \log(\hat{y_i}+1) )^2
    }
\end{equation}
La radice dell'MSLE.

% dire da qualche parte qualche altra valutaione: e.g. la computazione richiesta
% dai vari approcci

\subsubsection{Valutazione dei costi computazionali}

Il numero dei parametri allenabili dei vari modelli sono stati annotati insieme
al tempo richiesto per l'esecuzione di 10 epoche di train sulla piattaforma
\textit{Google Colab} abilitando l'accelerazione hardware GPU.

\subsection{Categoriche}
\textbf{Dire qualcosa sulle categoriche}

Inizialmente sono state valutate le performance di regressione dei prezzi a
partire dalle sole features categoriche in modo tale da stabilire un punto di
partenza per la successiva analisi delle features testuali \textit{name} e
\textit{item\_description}, più complesse e computazionalmente costose.

Un semplice modello composto da due layer Densi è stato sufficiente a
raggiungere le seguenti performance alla fine della decima epoca con l'aiuto di
un layer Dropout interposto ad essi per contrastare l'overfitting.

% todo: qui mettere le performance

% parlare di quanto tempo richiedono le 10 epoche.
% del numero dei parametri del modello.
% di quanti neuroni hanno i layers e perchè (magari fare una prova con vari
% valori bassi e aumentarli senza esagerare fino a raggiungere le performance sopra)
% parlare di qualche altro parametro? di training del modello?

\subsection{Rappresentazione del Testo}

\subsubsection{Bag of Words}
Il modello bag-of-words è una rappresentazione semplificata  di testi, dove
ciascuno di essi è rappresentato come una "borsa" (vettore) delle sue parole.
Quindi crea un vettore per ogni documento contenete il conteggio delle
occorrenze di ciascuna parola del vocabolario nel documento. Questa tecnica
ignora la grammatica e persino l'ordine delle parole, ma mantiene la
molteplicità \cite{manning_raghavan_schutze_2008}.
\subsubsection{Tf-Idf}\label{section-tfidf}
Tf-Idf (Term frequency-Inverse document frequency) \cite{manning_raghavan_schutze_2008} risolve uno dei principali problemi che si hanno con Bag of Words, poiché quest'ultima tecnica non considera relazioni tra i vari documenti. Infatti, Tf-Idf \eqref{eq:tf-idf} considera l'importanza di una parola per un documento in una collezione di documenti (corpus).
Il vettore costruito per ogni documento tramite questa tecnica considera per ciascun termine del documento la frequenza con cui il termine compare in questo documento (Term Frequency, TF \eqref{eq:tf}) rispetto alla frequenza del termine in tutti i documenti (Inverse Document Frequency, IDF \eqref{eq:idf}).

\begin{equation}
\label{eq:tf-idf}
   Tf-Idf_{t,d} = tf_{t,d} \cdot idf_t
\end{equation}\\
\begin{equation}
\label{eq:tf}
   tf_{t,d} = \frac{n_{t,d}}{number of terms in the document d} 
\end{equation}\\
\begin{equation}
\label{eq:idf}
   idf_{t} = \log \frac{number of documents}{number of documents with term t} 
\end{equation}
Dove in \eqref{eq:tf-idf}\eqref{eq:tf}\eqref{eq:idf} t indica il termine e d indica il documento.
\\
\subsubsection{Word Embeddings}
Le precedenti tecniche sono semplici da realizzare, robuste e funzionali per molti tasks. Tuttavia queste tecniche semplici sono al limite in molti compiti perché trattano le parole come unità atomiche, senza alcuna correlazione poiché vengono rappresentate come indici di un vocabolario. 
Con lo sviluppo delle tecniche di Machine Learning sono nati modelli in grado di apprendere rappresentazioni vettoriali di alta qualità di parole provenienti da set di dati composti da miliardi di parole e con vocabolari di milioni di parole. Queste nuove tecniche generano rappresentazioni vettoriali delle parole che rispettano la similarità tra le parole considerando più gradi di somiglianza, rappresentando la regolarità semantica e sintattica, permettendo anche operazioni algebriche \cite{mikolov2013efficient}.
\subsubsection{GLOVE}
GLOVE \cite{pennington2014glove} è un algoritmo di apprendimento non supervisionato che fornisce una rappresentazione vettoriale di ogni parola, considerando le statistiche globali di ricorrenza parola per parola (word-word ??????) da un corpus e le rappresentazioni ottenute mostrano sottostrutture lineari dello spazio vettoriale delle parole. Inoltre, gli ideatori di questa tecnica hanno reso disponibile i word embeddings pre-addestrati.
Nel dettaglio in questo lavoro sono stati utilizzati i seguenti embeddings:
\begin{itemize}
	\item Wikipedia 2014 + Gigaword 5: allenato su 6 miliardi di parole e con un vocabolario di 400 mila parole;
	\item Common Crawl: allenato su 840 miliardi di parole e con un vocabolario di 2.2 milioni di parole (il più grande disponibile di GLOVE);
\end{itemize}

\subsection{modello}
qualche spiegazione sulla definizione del modello.
perchè le lstm, cosa sono, paper, lstm per il testo.

performance gain delle lstm??
perchè non abbiamo provato le gru?

\subsection{Rappresentazione del Testo}
Descrizione dei vari approcci nella storia (quando sono stati creati)

deep learning e non

\subsubsection{clean o raw}
abbiamo valutato clean e raw; quali modelli beneficiano del cleanup, quali no?
perchè? il deep learning impara meglio se il testo è raw?

\begin{itemize}
    \item bag of word based models: count vectorizer - tf/idf
    \item embeddings: keras, glove pretrained. dire perchè non abbiamo provato
    qualcosa tipo word2vec?
    \item transformers con fine tuning
\end{itemize}



\section{Results and Evaluation}
The Results section is dedicated to presenting the actual results (i.e. measured and calculated quantities), not to discussing their meaning or interpretation. The results should be summarized using appropriate Tables and Figures (graphs or schematics). Every Figure and Table should have a legend that describes concisely what is contained or shown. Figure legends go below the figure, table legends above the table. Throughout the report, but especially in this section, pay attention to reporting numbers with an appropriate number of significant figures.
\subsection{Risultati clean}
\subsection{Risultati raw}

\subsubsection{clean o raw}
abbiamo valutato clean e raw; quali modelli beneficiano del cleanup, quali no?
perchè? il deep learning impara meglio se il testo è raw?


\begin{itemize}
    \item bag of word based models: count vectorizer - tf/idf
    \item embeddings: keras, glove pretrained. dire perchè non abbiamo provato
    qualcosa tipo word2vec?
    \item transformers con fine tuning
\end{itemize} 


\section{Discussion}
% The discussion section aims at interpreting the results in light of the
% project's objectives. The most important goal of this section is to interpret
% the results so that the reader is informed of the insight or answers that the
% results provide. This section should also present an evaluation of the
% particular approach taken by the group. For example: Based on the results, how
% could the experimental procedure be improved? What additional, future work may
% be warranted? What recommendations can be drawn?
\subsection{Modello Finale}
\subsubsection{Parametri di training}

\paragraph{Numero di neuroni:} Il numero di neuroni del modello finale dei vari
layer è mostrato nelle figure \ref{fig:modeltemplate} e \ref{fig:kerasModel}
allo stesso modo del rate di dropout.

\paragraph{Optimizer:} Come optimizer è stato utilizzato Adam impostando il
learning rate a 1e-4, in quanto il training è risultato più stabile per la
maggior parte degli approcci ad eccezione dell'Embedding Keras che ha dimostrato
un comportamento migliore con il learning rate di default (1e-3).

\paragraph{Batchsize:} La batchsize è stata impostata a 256.

\paragraph{Regolarizzazione:} Oltre ai layer Dropout è stato testata la
regolarizzazione L2, tuttavia non ha portato a evidenti miglioramenti ed è
quindi stata omessa. Inoltre è stato impostato l'earlystopping a monitoraggio
dell'errore sul validation impostando pazienza a 3 epoche e il ripristino
automatico dei pesi più performanti.

\subsubsection{Layer aggiuntivi per il testo}

Per quanto concerne i layer aggiuntivi in coda agli Embedding sono state
valutate layer ricorsivi (RNN) in quanto ottimi per l'elaborazione di dati
sequenziali, come il testo, essendo in grado di mantenere in "memoria"
informazioni sugli elementi precedenti \cite{liang2017text}. Nello specifico
sono stati valutati LSTM, GRU e Bidirectional LSTM. Queste ultime, permettono di
catturare relazioni sia con parole precedenti che con parole successive
\cite{schuster1997bidirectional}, ed essendosi sono dimostrate più performanti
compaiono nel modello finale.

Sono stati inoltre valutati layer convoluzionali a precedere i layer RNN; tuttavia
non avendo migliorato considerevolmente le performance sono stati omessi nel
modello finale.

L'aggiunta di un layer GlobalMaxPooling1D a seguire le Bi\_LSTM (figura
\ref{fig:kerasModel}) si è dimostrata efficace sia per rendere la dimensionalità
compatibile con i successivi layer Densi che nell'effettuare un
\textit{downsampling} riducendo la dimensionalità insieme ai costi
computazionali.

\subsubsection{Dimensione Word Embedding}
La dimensione del word embedding Keras, ovvero quello allenato insieme al resto
del modello, è stata impostata a 50.

\subsection{Embedding pretrainati}
In tabella \ref{tab:restable} si nota come i modelli word embedding
pre-allenati (GloVe) richiedano un numero di epoche maggiore per convergere.

\subsection{Pulizia del testo}

Generalmente il preprocessing completo del testo non ha prodotto evidenti
miglioramenti sulle performance in termini di errore rispetto al preprocessing
limitato, nonostante nel caso degli embedding pre-trainati il preprossing
completo permette una maggiore copertura delle parole del vocabolario: da 32\%
a 38\% per GloVe6B e da 42\% a 52\% per GloVe840B.

Per quanto riguarda gli embedding ciò può essere dovuto al fatto che essendo in grado di catturare relazioni semantiche tra le parole non
necessitano di un alto grado di preprossesing che rimuovendo le stopwords, ad
esempio, potrebbe perfino complicare l'apprendimento del significato della
frase.

Risulta invece inaspettato il risultato ottenuto con il preprocessing totale
applicato a Bag of Word e Tf-Idf: l'errore è infatti leggermente più alto rispetto
a quello con preprocessing limitato. Si deduce che operazioni come
rimozione di stopwords e lemmatizzazione non hanno offerto un contributo
rilevante ma anzi un leggero peggioramento. 

Si nota invece una generale diminuzione del numero di parametri (figura
\ref{fig:resParam}) che in alcuni casi si traduce in una riduzione del tempo
richiesto da una singola epoca di train (figura \ref{fig:resEtime}).

\subsection{Performance}
Il modello più performante dal punto di vista dell'errore di regressione risulta
l'embedding Keras allenato con il resto del modello, evidenziato in verde nella
tabella \ref{tab:restable}. Seguono gli approcci Bag of Word e Tf-Idf;
quest'ultimo tuttavia non sembra migliorare i risulatati del più semplice BoW.

Dal punto di vista computazionale invece, la durata di una singola epoca di Bag
of Word è notevolmente minore rispetto a quella dell'embedding Keras allenato,
tuttavia considerando il numero di epoche necessarie alla convergenza ci
accorgiamo che i tempi di training risultano simili: 5000s per BoW e 5530s per
Keras. È necessario ricordare che però Keras è stato allenato con un learning
rate maggiore che per BoW risultava invece in una scarsa stabilità.

Sono invece gli Embedding GloVe pretrainati ad essere i più veloci nella singola
epoca di training, senza tener conto del modello di riferimento tenente conto
delle sole feature non testuali, avendo il numero minore di
parametri allenabili come mostrato in figura \ref{fig:resParam}. Anch'esse
tuttavia richiedono numerose epoche per convergere. Non risultano esserci
notevoli differenza ne di errori ne di costi computazionali tra i due embedding
pretrainati GloVe6B e GloVe840B.

\subsection{BERT}

Non è stato possibile per motivazioni computazionali allenare questo
modello sulla totalità del dataset ma è stata utilizzata porzione del
dataset per il training (circa 200 mila istanze) completando solo 5 epoche.
Tuttavia considerando queste limitazioni i risultati sembrano promettenti.

È stata una piacevole sorpresa il risultato ottenuto da Bag of Word che pur
essendo l'approccio più semplice testato ottiene ottimi risultati sia in termini
di errore che di richieste computazionali rappresentandone un buon compresso.

\section{Conclusions}
% Conclusions should summarize the central points made in the Discussion section,
% reinforcing for the reader the value and implications of the work. If the
% results were not definitive, specific future work that may be needed can be
% (briefly) described. The conclusions should never contain ``surprises''.
% Therefore, any conclusions should be based on observations and data already
% discussed. It is considered extremely bad form to introduce new data in the
% conclusions.

Riassumendo dai risultati ottenuti si evince che il modello Keras permette di
ottenere ottime performance in cambio di costi computazionali accettabili,
perciò questa tecnica risulta consigliata nel caso in cui si abbia a
disposizione una piattaforma almeno comparabile con Google Colab.

Bag of Words è stata una piacevole sorpresa infatti pur essendo l'approccio più
semplice testato ottiene ottimi risultati ed è probabilmente attuabile anche
da chi ha particolari limitazioni computazionali.

È stato inoltre evidenziato che il pre-processing del testo effettuato
generalmente non giova alle performance e sarebbe interessante valutare approcci
alternativi a quello utilizzato nel progetto.

Nel caso si abbiano a disposizione grandi risorse computazionali sarebbe utile
valutare approfonditamente l'approccio basato su BERT che potrebbe portare a
performance ottimali.

Un altro approccio che varrebbe la pena testare sarebbe l'utilizzo di un embedding
pre-trainato con GloVe come punto di partenza per poi allenarne ulteriormente i
pesi magari ampliando le dimensioni dell'embedding in caso si abbiano risorse sufficienti.


\bibliographystyle{IEEEtran}
\bibliography{Project}

\end{document}